% Сборка из корня: lualatex -halt-on-error -interaction=nonstopmode -output-directory=SVD SVD/SVD_algorithm.tex
\documentclass[11pt,a4paper]{article}
\usepackage{fontspec}
\setmainfont{Noto Serif}
\usepackage[russian]{babel}
\usepackage[margin=21mm]{geometry}
\usepackage{amsmath,amssymb,microtype}
\setlength{\parindent}{0pt}
\setlength{\parskip}{4pt}
\setlength{\emergencystretch}{2em}
\newcommand{\R}{\mathbb R}
\newcommand{\Id}{\mathrm{Id}}
\newcommand{\T}{^{\mathsf T}}
\newcommand{\F}{_{\mathrm F}}
\DeclareMathOperator{\rank}{rank}
\begin{document}
\begin{center}
{\Large\bfseries SVD-решатель: последовательность алгоритма}\\[4pt]
{\small Краткое изложение в обозначениях \texttt{SVD/SVD\_solver.tex}}
\end{center}

\textbf{Задача и входные данные.}
Описывается один внутренний шаг ADP: текущий базис
$P\in\R^{m\times d}$, $PP\T=\Id_m$, и локальные величины
$g_j\in\R^m$, $I_j\in\R^p$, $U_j\in\R^{p\times d}$,
$N_j>0$, $j=1,\ldots,J$, фиксированы. Задаются $\lambda>0$,
предельный ранг $1\le r\le m$, допуски и предел числа внутренних итераций.
Решатель приближённо минимизирует
\[
 F(B)=\sum_{j=1}^J N_j\|I_j-U_jB\T g_j\|_2^2
       +\lambda\|B-P\|\F^2,
 \qquad \rank(B)\le r.
\]
Первое слагаемое согласует глобальную матрицу с локальными откликами;
второе удерживает её около текущего $P$ и делает условные линейные системы
положительно определёнными. Статистики $I_j,U_j$ имеют согласованную
нормировку; масса $N_j$ используется как вес, а не повторный множитель статистик.

\textbf{1. Инициализация и невязка.}
Положить $B_0=0$; список компонент пуст. Перед добавлением компоненты
$k=1,\ldots,r$ для $B=B_{k-1}$ вычислить
\[
 e_j=I_j-U_jB\T g_j,\qquad
 Q(B)=\sum_jN_j g_j(U_j\T e_j)\T+\lambda(P-B).
\]
Здесь $e_j$ --- невязки данных, а $Q=-\tfrac12\nabla F(B)$ --- матричное
направление уменьшения функционала. Ищем добавку $\sigma av\T$,
где $a\in\R^m$, $v\in\R^d$, $\|a\|_2=\|v\|_2=1$.
Одна такая добавка увеличивает ранг не более чем на единицу.

\textbf{2. Поиск пары направлений.}
Выбрать единичное начальное $v$: например, ведущий правый сингулярный
вектор $Q$ или случайный вектор. При неизменных $B,e_j,Q$ попеременно решать
\begin{align*}
 G_v\widetilde a&=Qv,&
 G_v&=\sum_jN_j\|U_jv\|_2^2g_jg_j\T+\lambda\Id_m,
 &a&=\widetilde a/\|\widetilde a\|_2,\\
 H_a\widetilde v&=Q\T a,&
 H_a&=\sum_jN_j(a\T g_j)^2U_j\T U_j+\lambda\Id_d,
 &v&=\widetilde v/\|\widetilde v\|_2.
\end{align*}
Первая система мала ($m\times m$). Вторую можно решать через действия
\[
 H_ax=\sum_jN_j(a\T g_j)^2U_j\T(U_jx)+\lambda x:
\]
плотную матрицу размера $d\times d$ строить необязательно.
Эквивалентно можно решать регуляризованную задачу наименьших квадратов
\[
 \min_x\;\sum_jN_j\|e_j-(a\T g_j)U_jx\|_2^2
             +\lambda\|x-(P-B)\T a\|_2^2.
\]
Попеременные шаги улучшают пару с учётом различной чувствительности
локальных операторов. Остановить поиск при стабилизации направлений
(например, малости $1-|v_{t+1}\T v_t|$) или достижении лимита;
контролировать также невязки систем. Нулевые $\widetilde a$ или
$\widetilde v$ не нормировать: прекратить попытку либо выбрать другое начало.

\textbf{3. Точный масштаб и принятие компоненты.}
Для найденной пары вычислить
\[
 R_B(a,v)=a\T Q(B)v,\qquad
 D(a,v)=\sum_jN_j(a\T g_j)^2\|U_jv\|_2^2+\lambda,
 \qquad \sigma_\star=\frac{R_B(a,v)}{D(a,v)}.
\]
Поскольку
$F(B+\sigma av\T)-F(B)=D\sigma^2-2R_B\sigma$,
это точный минимум по масштабу; ожидаемое уменьшение равно
$\gamma=R_B^2/D$. Проверить фактическое уменьшение исходного $F$,
затем добавить $\sigma_\star av\T$. Знак масштаба можно перенести в $a$.
При слишком малом улучшении завершить построение либо проверить другое начало.

\textbf{4. Сжатие и совместная переоценка.}
Записать накопленную матрицу как $LR\T$, выполнить сокращённые
QR-разложения $L=Q_LR_L$, $R=Q_RR_R$ и малое SVD
$R_LR_R\T=\widehat A\Sigma\widehat V\T$. Тогда
\[
 B=A\Sigma V\T,\qquad A=Q_L\widehat A,\quad V=Q_R\widehat V,
 \qquad A\T A=V\T V=\Id_s,
\]
где $s\le k$ --- число ненулевых компонент. Это устраняет зависимые
факторы без изменения матрицы. Удаление ненулевых малых сингулярных
значений уже меняет $B$ и требует отдельной проверки $F$.
Для фиксированных столбцов $a_\ell,v_\ell$ совместно уточнить масштабы:
\begin{align*}
 M_{\ell t}&=\sum_jN_j(a_\ell\T g_j)(a_t\T g_j)
        (U_jv_\ell)\T(U_jv_t)+\lambda\delta_{\ell t},\\
 h_\ell&=\sum_jN_j(a_\ell\T g_j)I_j\T U_jv_\ell
                 +\lambda a_\ell\T Pv_\ell,\qquad M\sigma=h.
\end{align*}
Положить $B_k=\sum_{\ell=1}^s\sigma_\ell a_\ell v_\ell\T$.
Малая система учитывает взаимодействие всех уже выбранных компонент
и даёт точный минимум в их линейной оболочке; при $\lambda>0$ она
положительно определена. Коэффициенты $\sigma_\ell$ могут иметь любой знак.

\textbf{5. Остановка и результат.}
Повторять шаги 1--4 до $r$ добавок или малости относительного улучшения
$\gamma/\max(1,F(B_{k-1}))$; сохранять итоговую $B$, значение $F$ и причину
остановки. Попеременный поиск не гарантирует глобального минимума:
малое улучшение найденной пары не исключает более полезных направлений.
Для следующего внешнего шага ADP извлекают ортонормированный базис строк $B$.
При $\rank(B)<m$ недостающие направления дополняют из $P$ после проекции
на дополнение к найденному пространству; эти направления не восстановлены
самой $B$. Если $B=T\widehat P$ имеет полный ранг строк, то переход
к $\widehat P\widehat P\T=\Id_m$ сохраняет предсказания при замене
$g_j\mapsto T\T g_j$. Ортонормирование не обязано сохранять штраф в $F$.

\textbf{Вариант малоранговой поправки.}
Для ограничения $\rank(\Delta)\le r$, $B=P+\Delta$, применить ту же схему
к $\Delta$, начиная с $\Delta_0=0$, с заменами
$I_j\mapsto I_j-U_jP\T g_j$ и регуляризационной цели $P\mapsto0$.
Вернуть $B=P+\Delta$: малый ранг здесь относится к изменению базиса,
а полный ранг $B$ необходимо проверять отдельно.

\textbf{Отличие от текущей реализации (\texttt{SVD\_cur.tex}).}
При стандартном штрафе Фробениуса математическая схема совпадает;
\texttt{SVD.py} конкретизирует её вычислительные шаги. Функция
\texttt{solve} сама оценивает $g_j$ на входном базисе и фиксирует их
до завершения жадного построения, а затем повторно оценивает на итоговом
базисе. Начальное $v$ выбирается по ведущей компоненте
$Q-(QV)V\T$, то есть вне уже найденного правого пространства;
дальнейший поиск не ограничивается этим дополнением. Условные системы
решаются Холецким или LSMR с проверкой нормальной невязки и, при
необходимости, уточнением; используются кэш $U_jV$, проверки численного
ранга и отсутствия роста функционала. Реализация требует $r<m$,
допускает $\lambda=0$ и возвращает ортонормированный базис с диагностикой
(отдельный \texttt{solve\_fixed\_coefficients} возвращает сырую $B$).
Опциональная спектральная метрика меняет регуляризационный штраф и
реализуется преобразованием координат; в изложенной выше базовой схеме
она не включена.
\end{document}
